\section{Constructive matrix proofs of rigidity and extremizer classification}
\label{sec:matrix-rigidity}

The equality and stability results in this section concern actual complex
density matrices and actual qubit unitaries.  In particular, no unproved
rotation-to-unitary lifting assertion is used.  Write
\(\tau(\sigma)=\|T(\sigma)\|_F^2\), and use the trace norm
\(\|X\|_1=\sum_j|\lambda_j(X)|\) for Hermitian \(X\).
The usual trace distance between states is \(\frac12\|\rho-\sigma\|_1\).

\subsection{A complete radial normal form, including degenerate cases}

For a canonical state, Bob's marginal is \(I/2\).  First diagonalize Alice's
marginal by a qubit unitary, putting her Bloch vector at \(c e_z\), where
\(c\geq0\).  In Alice blocks write
\(\sigma=\left(\begin{smallmatrix}A&C\\C^\dagger&D\end{smallmatrix}\right)\).
Then \(A+D=I/2\).  A second qubit unitary, acting only on Bob, diagonalizes
the Hermitian matrix \(H=A-D\), with its eigenvalues in decreasing order.
If their difference is \(t\geq0\), the resulting state is exactly
\begin{equation}
 M(c,t,z,w,v)=
 \begin{pmatrix}
 a_+&0&z&w\\
 0&a_-&v&-z\\
 \overline z&\overline v&d_-&0\\
 \overline w&-\overline z&0&d_+
 \end{pmatrix},\qquad
 a_\pm=\frac{1+c\pm t}{4},\quad
 d_\pm=\frac{1-c\pm t}{4}.
 \label{eq:rigidity-normal-form}
\end{equation}
Positivity implies \(0\leq t\leq1-c\).  Both spectral diagonalizations
remain valid for repeated eigenvalues and singular density matrices.
The trace-zero condition on \(C\), which forces its diagonal entries to
be \(z,-z\), follows from the vanishing off-diagonal entry of Alice's
marginal.

There is also a version with a prescribed radial axis.  If Alice's Bloch
vector is already \(c e_z\), diagonalize only \(H=A-D\).  Direct Pauli
expansion gives
\begin{equation}
 \operatorname{BlochSq}(H)=\|T^{\mathsf T}e_z\|^2=t^2.
\end{equation}
Thus the prescribed \(t\) is preserved.  This statement holds even when
\(c=0\): the marked axis is still \(e_z\), and Alice's unitary is exactly
the identity.  Diagonalizing a maximally mixed Alice marginal alone would
not justify preservation of an arbitrary marked axis.

In the normal form,
\begin{equation}
 \tau=t^2+16|z|^2+8|w|^2+8|v|^2.
 \label{eq:normal-form-correlation}
\end{equation}
The sharp geometric inequality \(\tau\leq1-c^2+2t^2\) has the following
complete equality conditions:
\begin{align}
 t z&=0,\label{eq:rigidity-equality-one}\\
 |z|^2+|w|^2&=a_+d_+,\label{eq:rigidity-equality-two}\\
 |z|^2+|v|^2&=a_-d_-,\label{eq:rigidity-equality-three}\\
 z\overline v&=w\overline z.\label{eq:rigidity-equality-four}
\end{align}
These are necessary and sufficient for every positive matrix
\eqref{eq:rigidity-normal-form}.  For \(t>0\), they reduce to
\begin{equation}
 z=0,\qquad |w|^2=a_+d_+,\qquad |v|^2=a_-d_-.
 \label{eq:positive-t-equality}
\end{equation}
Thus all extremizers in this basis are the explicit two-block states
\eqref{eq:positive-t-equality}, including arbitrary admissible complex
phases.  At \(t=0\), the correct condition is instead
\begin{equation}
 C C^\dagger=\frac{1-c^2}{16}I.
 \label{eq:zero-t-equality}
\end{equation}
In particular, setting \(z=0\) at \(t=0\) would discard legitimate
extremizers.  Equation \eqref{eq:zero-t-equality} includes \(c=1\), when
\(C=0\).

For necessity, the geometric spin-flip proof gives the exact identity
\[
 4\operatorname{Re}\operatorname{tr}(M\widetilde M)
 =1-c^2+2t^2-\tau.
\]
Both \(M\) and \(\widetilde M\) are positive.  Zero trace therefore implies
\(M\widetilde M=0\): conjugate by the positive square root of
\(\widetilde M\), use that a positive matrix with zero trace is zero, and
then use the kernel of that square root.  Four entries of the zero product
give \eqref{eq:rigidity-equality-one}--\eqref{eq:rigidity-equality-four}.
Conversely, the two norm identities and
\eqref{eq:normal-form-correlation} give equality immediately.  This proof
never divides by an eigenvalue or a diagonal entry.

\subsection{Actual matrix stability for the canonical envelope}

Put \(d=1-c\), and define the canonical maximizing reference
\begin{equation}
 M_{\rm ob}(c)=\frac12
 \begin{pmatrix}
 1&0&0&\sqrt d\\
 0&c&0&0\\
 0&0&0&0\\
 \sqrt d&0&0&d
 \end{pmatrix}.
 \label{eq:aligned-obese}
\end{equation}
This is a Bob Pauli-X conjugate of the standard maximally obese state.
For every canonical state \(\sigma\), let
\[
 c=|a|,\qquad
 \delta=(1-c)(3-c)-\tau(\sigma)\geq0.
\]
If \(0\leq c<1\), there exist genuine qubit unitaries \(U,V\) such that,
with \(\sigma'=(U\otimes V)\sigma(U\otimes V)^\dagger\),
\begin{equation}
 \|\sigma'-M_{\rm ob}(c)\|_1^2
 \leq4\|\sigma'-M_{\rm ob}(c)\|_{\rm HS}^2
 \leq \delta\left(\frac12+\frac{2+c}{1-c+t}\right)
 \leq\frac{5+c}{2(1-c)}\,\delta,
 \label{eq:actual-matrix-rigidity}
\end{equation}
where \(t\) is the radial parameter of the normal form chosen above.
The trace-norm estimate can be capped by the diameter of the state space:
\begin{equation}
 \|\sigma'-M_{\rm ob}(c)\|_1^2
 \leq \min\left\{4,\frac{5+c}{2(1-c)}\delta\right\}.
 \label{eq:capped-matrix-rigidity}
\end{equation}
Consequently the usual trace distance is at most
\(\min\{1,\frac12\sqrt{(5+c)\delta/[2(1-c)]}\}\).

Here is the matrix estimate behind the theorem.  An Alice diagonal phase
unitary makes \(w\) real and nonnegative; the zero-coherence case is
handled by choosing phase one.  Set \(y=d-t\), \(w_0=\sqrt d/2\),
\(Z=|z|^2\), and \(J=|v|^2\).  Direct expansion of the actual entries gives
\begin{align}
 \|M-M_{\rm ob}(c)\|_{\rm HS}^2
   &=\frac{y^2}{4}+4Z+2J+2(w-w_0)^2,\\
 \delta&=d^2-t^2+8(w_0^2-w^2)-16Z-8J.
\end{align}
Positive principal minors give
\[
 Z\leq\frac{(1+c+t)y}{16},\qquad
 J\leq\frac{(1+c-t)y}{16},\qquad
 0\leq w\leq w_0.
\]
Using \((w-w_0)^2\leq w_0^2-w^2\) yields
\[
 4\|M-M_{\rm ob}(c)\|_{\rm HS}^2
 \leq\delta+y(3+3c-t).
\]
The geometric inequality gives
\(\delta\geq2y(d+t)\).  Substitution proves the middle estimate in
\eqref{eq:actual-matrix-rigidity}; monotonicity in \(t\geq0\) gives the last.
The first estimate follows from the spectral theorem and Cauchy--Schwarz
for four eigenvalues.  The cap follows from positivity and trace one of
both density matrices.  These spectral and matrix identities are proved
in Lean, rather than supplied as assumptions to scalar algebra.

Setting \(\delta=0\) proves the complete canonical equality classification:
\begin{equation}
 \tau(\sigma)=(1-c)(3-c)
 \quad\Longleftrightarrow\quad
 \exists U,V:\ (U\otimes V)\sigma(U\otimes V)^\dagger=M_{\rm ob}(c).
 \label{eq:canonical-LU-classification}
\end{equation}
The endpoint \(c=1\) is proved separately: the normal form has \(t=0\),
and positivity forces all three coherences to vanish.  No estimate
containing \(1/(1-c)\) is used at that endpoint.

\subsection{The complete global maximizing orbits}

For \(0\leq c<2/3\), put
\[
 D(c)=2-2c-c^2,\qquad r_*(c)=\frac{c(1-c)}{D(c)},\qquad d=1-c.
\]
Define the density matrix
\begin{equation}
 \rho_*(c)=\frac12
 \begin{pmatrix}
 1+r_*&0&0&\sqrt{d(1-r_*^2)}\\
 0&c(1-r_*)&0&0\\
 0&0&0&0\\
 \sqrt{d(1-r_*^2)}&0&0&d(1-r_*)
 \end{pmatrix}.
 \label{eq:global-optimal-normal-form}
\end{equation}
Every state with actual steering-centre length \(c\) satisfies
\begin{equation}
 \tau(\rho)=G(c)
 \quad\Longleftrightarrow\quad
 \exists U,V:\ (U\otimes V)\rho(U\otimes V)^\dagger=\rho_*(c).
 \label{eq:global-LU-classification}
\end{equation}
The theorem assumes no rank condition.  If equality holds, \(G(c)>1\)
forces Bob's marginal to be invertible, because a pure Bob marginal gives
\(\tau\leq1\).  The global deficit identity then forces the canonical
representative to saturate its canonical envelope.  Apply
\eqref{eq:canonical-LU-classification} and conjugate the positive
reconstruction filter by the same Bob unitary.

The remaining filter orientation is fixed by an exact identity.  For the
canonical reference \eqref{eq:aligned-obese} and any normalized complex
filter with effect \(F^\dagger F=I+b\cdot\sigma\),
\begin{equation}
 G(c)-\tau(FM_{\rm ob}(c)F^\dagger)
 =(2-3c)(b_x^2+b_y^2)+D(c)(b_z-r_*)^2.
 \label{eq:obese-directional-completed-square}
\end{equation}
The displayed notation means filtering on Bob.  Both coefficients are
strictly positive for \(c<2/3\), so every maximizing effect is exactly
\(I+r_*\sigma_z\).  The reconstruction filter is positive, hence uniqueness
of the positive square root gives
\(F=\operatorname{diag}(\sqrt{1+r_*},\sqrt{1-r_*})\), which proves
\eqref{eq:global-optimal-normal-form}.  The effect argument covers all
complex normalized filters, including singular ones; no unproved polar
decomposition is needed.

For \(2/3\leq c<1\), the supremum \(G(c)=1\) is not attained, as proved by
the strict flat-branch bound for both marginal-rank cases.  At \(c=1\),
the only maximizers are pure product states.  Indeed an invertible Bob
marginal would give the strict joint bound \(\tau<1\); with singular Bob,
\(\tau=1\) forces Alice's marginal to be pure too.  Ordered qubit
diagonalization puts both marginals at \(|0\rangle\langle0|\), and the
positive support identity then forces the entire joint state to be
\(|00\rangle\langle00|\).

\subsection{Formal coverage of the matrix statements}

The new modules use the same pinned Lean 4.19.0 and mathlib revision as
the original development.  The principal declarations are
\path{canonical_aligned_normal_form},
\path{physical_bob_radial_normal_form},
\path{physical_radial_equality_iff},
\path{canonical_trace_rigidity_capped},
\path{canonical_obese_equality_iff},
\path{alignedObese_optimal_effect},
\path{global_steering_equality_iff}, and
\path{global_endpoint_equality_iff}.
Their hypotheses are positivity, normalization, and the stated measured
coordinates; actual unitary witnesses, matrix distances, and singular
support are derived inside Lean.  The marked-axis theorem uses \(e_z\)
explicitly, including at zero centre.  The global maximizing-orbit
classification has no coordinate or local-unitary lifting hypothesis.

